\section{Deformation inside a division algebra}\label{sec:division}

The rank bound from Proposition~\ref{prop:rectangular} is close to full
rank.  We now deform the test matrices so that every possible rank is a
multiple of their block size.  The same deformation will also turn a
defined nonzero value of a rational formula into an invertible value over
an extension field.
The construction is a form of the generic cyclic, or symbol,
division algebra. Explicit split representations and block-rank
divisibility appear in the constructive regularity methods of
Ivanyos, Qiao, and Subrahmanyam~\citep[Sections~4.3 and~5.2 of the
preprint]{IQS17}; see also their later regularity
algorithm~\citep[Section~4]{IQS18}.
Division-algebra witnesses and invertible hits are also used in
rational-formula black-box testing by Arvind, Chatterjee, and
Mukhopadhyay~\citep{ACM22,ACM25}.
Here the needed deformation and rank-divisibility statement have
a direct proof using two explicit matrices. This matrix embedding
is distinct from the rational multiplication representation in
Section~\ref{sec:generator}.

\begin{lemma}[A matrix lift into a division algebra]\label{lem:division-lift}
Let $d\geq 1$, let $K$ be a field of characteristic zero, and let
$\omega\in K$ be a primitive $d$th root of unity.  On the standard basis
$e_0,\ldots,e_{d-1}$ of $K^d$, with subscripts taken modulo $d$, set
\[
  De_j=\omega^j e_j,
  \qquad Se_j=e_{j+1}.
\]
For algebraically independent indeterminates $p,r$ over $K$, there is a
division algebra
\[
  \mathcal D\subseteq\Mat_d\bigl(K(p,r)\bigr)
\]
containing, for every $R\in\Mat_d(K)$, the matrix
\begin{equation}\label{eq:division-lift}
  R^*=\sum_{a,b=0}^{d-1}c_{ab}p^a r^bD^aS^b,
  \qquad
  c_{ab}=\frac1d\sum_{j=0}^{d-1}\omega^{-aj}R_{j,j-b}.
\end{equation}
The matrix $R^*$ specializes to $R$ at $p=r=1$.
Moreover, $\mathcal D$ has dimension $d^2$ over the central subfield
$K(p^d,r^d)$, and every matrix with blocks in $\mathcal D$ has rank over
$K(p,r)$ divisible by $d$.
\end{lemma}

\begin{proof}
For each $b$, the only potentially nonzero entries of $D^aS^b$ have the
form
\[
  (D^aS^b)_{j,j-b}=\omega^{aj}.
\]
The identity
\[
  \sum_{j=0}^{d-1}\omega^{(a-a')j}
  =\begin{cases}d,&a=a',\\0,&a\ne a',\end{cases}
  \qquad 0\leq a,a'<d,
\]
therefore gives Fourier inversion separately on each of these cyclic
diagonals.  In particular, the $d^2$ matrices $D^aS^b$ form a basis of
$\Mat_d(K)$, with the coefficients displayed in
\eqref{eq:division-lift}.  This proves the asserted specialization.

Put $X=pD$ and $Y=rS$, and let $\mathcal A$ be the unital
$K$-algebra they generate inside $\Mat_d(K(p,r))$.  Since
$DS=\omega SD$, we have
\begin{equation}\label{eq:division-monomials}
  XY=\omega YX,
  \qquad
  (X^iY^j)(X^kY^\ell)
    =\omega^{-jk}X^{i+k}Y^{j+\ell}
  \quad(i,j,k,\ell\geq0).
\end{equation}
Thus the ordered monomials $X^iY^j$ span $\mathcal A$ over $K$.
They are also linearly independent: in any finite relation, the
coefficient of the scalar monomial $p^i r^j$ is the corresponding
scalar coefficient times the invertible constant matrix $D^iS^j$.
Algebraic independence of $p,r$ forces each such coefficient to vanish.

Order the exponent pairs lexicographically.  For two nonzero elements
of $\mathcal A$, the largest exponent in their product is the sum of
their largest exponents, by \eqref{eq:division-monomials}.  Its
coefficient is the product of two nonzero coefficients and a nonzero
power of $\omega$.  Hence $\mathcal A$ has no zero divisors.

The elements
\[
  X^d=p^dI_d,
  \qquad Y^d=r^dI_d
\]
are scalar matrices.  We may consequently invert every nonzero
element of the central polynomial ring $K[p^d,r^d]I_d$ and obtain a
subalgebra $\mathcal D$ of $\Mat_d(K(p,r))$.  Every element of this
localization can be written $a/z$ with $a\in\mathcal A$ and
$0\ne z\in K[p^d,r^d]$.  Products of nonzero such fractions remain
nonzero, so $\mathcal D$ still has no zero divisors.

Let $F_0=K(p^d,r^d)$.  Reducing the exponents of $X,Y$ modulo $d$
shows that
\begin{equation}\label{eq:division-central-basis}
  \{X^aY^b:0\leq a,b<d\}
\end{equation}
spans $\mathcal D$ over $F_0$.  To see independence over $F_0$,
clear a common denominator in a proposed relation.  Expanding the
resulting coefficients in $K[p^d,r^d]$ gives a relation among the
ordered monomials over $K$; distinct pairs $(a,b)$ in
\eqref{eq:division-central-basis} have distinct exponent residues
modulo $d$.  All its coefficients must therefore vanish.
It follows that $\dim_{F_0}\mathcal D=d^2$.

For a nonzero $z\in\mathcal D$, left and right multiplication by $z$
are injective $F_0$-linear maps from this finite-dimensional space to
itself, and hence are surjective.  There exist $b,c\in\mathcal D$
with $zb=1$ and $cz=1$, and associativity gives
$c=c(zb)=(cz)b=b$.  Thus $z$ has a two-sided inverse in
$\mathcal D$.  This proves that $\mathcal D$ is a division algebra.
Since $p^ar^bD^aS^b=X^aY^b$, every lift in
\eqref{eq:division-lift} belongs already to $\mathcal A$.

Finally, consider an arbitrary block matrix with entries in
$\mathcal D$.  If it has a nonzero block, move that block to the
first position by block row and column permutations.  Writing the
result as
\[
  \begin{pmatrix}a&b\\c&E\end{pmatrix},\qquad a\ne0,
\]
block elimination gives
\[
  \begin{pmatrix}1&0\\-c&I\end{pmatrix}
  \begin{pmatrix}a^{-1}&0\\0&I\end{pmatrix}
  \begin{pmatrix}a&b\\c&E\end{pmatrix}
  \begin{pmatrix}1&-a^{-1}b\\0&I\end{pmatrix}
  =
  \begin{pmatrix}1&0\\0&E-ca^{-1}b\end{pmatrix}.
\]
All blocks and all these invertible operations remain in
$\mathcal D$.  Iterating reduces the matrix to identity blocks and
zero blocks.  Each identity block is an ordinary $d\times d$
identity matrix, and the operations are invertible over $K(p,r)$
as well.  Its ordinary rank is therefore $d$ times the number of
identity blocks, as required.
\end{proof}

The scalar fields in this lemma have different roles.  Finite
dimensionality and the division property are proved over the central
field $K(p^d,r^d)$.  Ordinary matrix rank is measured over the larger
field $K(p,r)$ in the given embedding.  The elimination proof uses
only operations in $\mathcal D$; it does not enlarge that algebra by
arbitrary scalars from $K(p,r)$.

\begin{proposition}[Full-rank deformed pencils]\label{prop:pencil-deformation}
Fix a primitive $M$th root of unity $\omega\in\C$, and lift the
padded matrices $R_0(t),\ldots,R_n(t)$ by
Lemma~\ref{lem:division-lift} with $d=M$ and
$K=\Q(\omega)(t)$.  Let
\[
  L(x)=A_0+\sum_{i=1}^n A_i x_i,
  \qquad A_i\in\Mat_q(\Q),\qquad 1\leq q\leq B,
\]
be an affine pencil admitting an invertible matrix evaluation over
some extension field of $\Q$.  Then
\begin{equation}\label{eq:deformed-pencil-determinant}
  P_{L,\omega}(t,p,r)
   :=\det\!\left(\sum_{i=0}^n A_i\otimes R_i^*(t,p,r)\right)
   \ne0.
\end{equation}
This is a polynomial in $\Q(\omega)[t,p,r]$ of total degree at most
$3BM^2$.  In particular $\det R_0^*$ is a nonzero polynomial.
\end{proposition}

\begin{proof}
Proposition~\ref{prop:rectangular} gives
\[
  \rank_{\Q(t)}\left(\sum_{i=0}^n A_i\otimes R_i(t)\right)
      \geq qN>(q-1)M.
\]
Choose a nonzero minor of order greater than $(q-1)M$.
It remains nonzero over $\Q(\omega)(t)$.  The corresponding minor
of the lifted matrix specializes to it at $p=r=1$, so the lifted
minor is not identically zero.  The lifted matrix consequently has
rank greater than $(q-1)M$ over $\Q(\omega)(t,p,r)$.

Its $M\times M$ blocks belong to the division algebra of
Lemma~\ref{lem:division-lift}: each is a rational linear combination
of the $R_i^*$.  Its rank is thus a multiple of $M$.  As the whole
matrix has order $qM$, its rank must be $qM$, proving
\eqref{eq:deformed-pencil-determinant}.

The entries of $R_i(t)$ lie in $\mathbb Z[t]$ and have degree at
most $N-1$.  The Fourier coefficients in
\eqref{eq:division-lift} introduce only the constant denominator
$M$, while the factors $p^ar^b$ have total degree at most $2(M-1)$.
Every lifted entry is therefore a polynomial of total degree at
most $N-1+2(M-1)<3M$.  A determinant of order $qM$ has total degree
at most $3qM^2\leq3BM^2$.

For the last assertion, apply the result to the order-one constant
pencil $L=1$, whose coefficient $A_0$ is $1$ and whose other
coefficients vanish.  Its determinant in
\eqref{eq:deformed-pencil-determinant} is exactly $\det R_0^*$.
\end{proof}

At any specialization at which $R_0^*$ is invertible, set
$X_i=(R_0^*)^{-1}R_i^*$ for $1\leq i\leq n$.  The identity
\begin{equation}\label{eq:division-left-normalization}
  A_0\otimes I_M+\sum_{i=1}^n A_i\otimes X_i
  =
  \bigl(I_q\otimes(R_0^*)^{-1}\bigr)
  \left(\sum_{i=0}^n A_i\otimes R_i^*\right)
\end{equation}
shows that simultaneous nonvanishing of the determinant constraints
in Proposition~\ref{prop:pencil-deformation}, including the
constant-pencil constraint, makes the corresponding affine pencils
invertible on this normalized tuple.  The order of multiplication
in \eqref{eq:division-left-normalization} is part of the identity.

\begin{lemma}[An invertible witness for a nonzero formula]
\label{lem:invertible-witness}
Every admissible nonzero rational formula over $\Q$ has a defined
evaluation at matrices over an extension field of $\Q$ for which
its value is invertible.
\end{lemma}

\begin{proof}
By the definition of nonzero, there is a tuple
$W=(W_1,\ldots,W_n)$ of rational matrices of some common dimension
$d\geq1$ at which the given formula is defined and has nonzero
value.  Choose a primitive $d$th root $\omega$, put $K=\Q(\omega)$,
and lift each $W_i$ to $W_i^*$ by
Lemma~\ref{lem:division-lift}.

We first verify the domain of the actual formula tree independently
of the division-algebra argument.  Inside $K(p,r)$ consider the
commutative ring
\[
  \mathcal O=
  \left\{\frac{f}{g}: f,g\in K[p,r],\ g(1,1)\ne0\right\}.
\]
Evaluation at $(1,1)$ is a ring homomorphism
$\varepsilon:\mathcal O\to K$.  An element of $\mathcal O$ whose
image under $\varepsilon$ is nonzero is a unit in $\mathcal O$:
if it is written $f/g$, then $f(1,1)\ne0$ and its inverse is $g/f$.

We claim, by induction up the original tree, that each subtree
evaluated at the lifted inputs gives a matrix over $\mathcal O$
whose image under $\varepsilon$ is its value at $W$.
This holds for leaves by the specialization property of the lift,
and is preserved by addition and ordered multiplication.  At an
inverse gate, let $G\in\Mat_d(\mathcal O)$ be the already
constructed value of the child.  Since the original evaluation is
defined, $\varepsilon(G)$ is invertible.  Consequently
\[
  \varepsilon(\det G)=\det\varepsilon(G)\ne0,
\]
so $\det G$ is a unit of $\mathcal O$.  The formula
\[
  G^{-1}=\frac{\operatorname{adj}(G)}{\det G}
\]
then constructs its inverse in $\Mat_d(\mathcal O)$ and gives the
correct specialized inverse.  This proves the claim at every
inverse gate, regardless of its nesting depth.

The lifted formula is therefore defined over $K(p,r)$.  Its final
matrix is nonzero, because one of its entries specializes to a
nonzero entry of the original final value.

Now use the division algebra $\mathcal D$ from
Lemma~\ref{lem:division-lift}.  Every lifted input and every
rational scalar constant belongs to $\mathcal D$.  A second
induction shows that every intermediate value belongs to
$\mathcal D$: sums and products stay in it, and an operand at an
inverse gate is already known to be invertible as an ambient
matrix, hence is nonzero in $\mathcal D$.  Its inverse in
$\mathcal D$ equals its ordinary matrix inverse by uniqueness.
The final nonzero value is thus a nonzero element of a division
algebra, and is invertible as a matrix over $K(p,r)$.
\end{proof}

Only those entries proved to lie in $\mathcal O$ are specialized in
this argument.  No specialization map on the whole division algebra
is needed.  In particular, the inverse of the final value may have
a pole at $(1,1)$ when the original nonzero value is singular; its
invertibility is asserted only at the lifted tuple over the
extension field.  If $d=1$, the lift leaves every input scalar
unchanged, and a nonzero scalar final value is already invertible.
The witness dimension $d$ in Lemma~\ref{lem:invertible-witness} is
purely existential and need not be bounded.  The matrices used in
the construction instead have the predetermined block size $M$.
